\section{Related Work}

\subsection{LLM Repair Loops with Execution Feedback}

The execution-feedback repair paradigm was established by Self-Debug
\citep{Chen2023Teaching}, which demonstrates that providing LLMs with their
own program output --- captured stdout, stderr, and exception traces --- enables
iterative repair that substantially improves pass@k over single-attempt
generation. Self-Debug operates on execution-only feedback, treating the
binary pass/fail signal augmented with error output as sufficient for repair.
This becomes our \textbf{Condition A baseline}: the de-facto standard that we
test against.

Reflexion \citep{Shinn2023Reflexion} extends the repair paradigm by having the
LLM generate verbal \emph{reflections} on its failures before attempting repair.
Verbal reflection is unstructured: the LLM authors its own diagnostic narrative.
Our work tests a complementary hypothesis --- that \emph{external} structured
diagnostics (mypy output) provide signal beyond what the LLM can self-generate
or extract from execution output.

CodeT \citep{Chen2022CodeT} and AlphaCode \citep{Li2022Competition} use
execution feedback at generation time (test case filtering, tournament selection)
rather than in iterative repair loops. These are orthogonal to our setting: we
fix the generation strategy and vary the repair feedback type.

\textbf{Gap we address:} Prior repair loop work treats feedback channel selection
as an engineering choice, not an experimental variable. Self-Debug, Reflexion,
and their successors do not ablate execution vs.\ execution+static-analysis.
We provide the first controlled comparison.

\subsection{Static Analysis in Code Generation and Repair}

Static analysis tools have been integrated into LLM code generation systems,
but without controlled ablation of their marginal contribution. SWE-agent
\citep{Yang2024SWEagent} uses a rich tool-use environment including bash
execution, file editing, and code analysis, embedded in an agent control loop
(ACI). The agent can invoke linters, but the ACI design makes it difficult to
isolate the contribution of any single feedback channel. We specifically avoid
this bundling: our experiment holds all factors constant except the feedback type.

LLM-based program repair (APR) surveys \citep{Xia2023Automated} document that
execution-based feedback is the dominant repair signal in the literature.
Type-system feedback appears in some systems, but controlled comparison against
execution-only baselines on Python benchmarks is absent.

\textbf{Gap we address:} Static analysis is routinely bundled with execution
feedback in production systems, but its isolated marginal contribution has not
been measured on standard Python code generation benchmarks.

\subsection{Evaluation Benchmarks}

EvalPlus \citep{Liu2023EvalPlus} introduces HumanEval+ and MBPP+ --- enhanced
versions of HumanEval \citep{Chen2021Evaluating} and MBPP \citep{Austin2021Program}
with substantially more test cases. We use HumanEval+ and MBPP+ as our primary
evaluation benchmarks. A contribution of our work is a new structural observation:
HumanEval+ and MBPP+ have dramatically different mypy error profiles in
GPT-4o-mini failing solutions (70\% vs 0\%). Prior work does not report mypy
error rates by benchmark, leaving this divergence unobserved.

\subsection{Formal Verification for Code Correctness}

Recent work explores whether LLMs can generate formal specifications that
constraint solvers then check. Our h-z1 sub-hypothesis tests whether
LLM-generated Z3 specs can provide useful counterexample feedback for repair,
finding that while spec validity is high (91.1\%), the counterexample activation
rate is low (16.1\%) on our arithmetic HumanEval+ subset. The neural-symbolic
program synthesis literature (e.g., Rosette \citep{Torlak2014Growing}) uses SMT
solvers for specification-guided synthesis; our setting uses LLM-generated
specifications, trading completeness for automation.
